\subsection{CTR: fuerza de predicción e inversión de la transición}
\label{sec:ctr-finite-reductions}
El §214 distingue el coste de acertar \(n\) bits de la ventaja de distinguir
un transcript completo. Para evaluar predicción usamos un suceso unilateral.
Para la transición usamos la terna desenmascarada de la última actualización.
Las dos derivaciones comparten la reducción por tramos de
la sección~\ref{sec:ctr-prp}. Escribimos \(M=2^{128}\) y
\((M)_r=M(M-1)\cdots(M-r+1)\). La fuerza concreta de AES figura como
\(\alpha_r(T)\), la ventaja PRP con clave uniforme, hasta \(r\)
consultas y tiempo total \(T\) incluida la simulación.

\paragraph{1. Predicción de varios bloques (.5--.8).}
Fijado el prefijo sustituido, un tramo consulta \(r\in\{3,4\}\)
puntos distintos. Una función aleatoria da \(r\) palabras independientes.
Una permutación aleatoria da esa misma ley condicionada a que sean distintas.
Para cualquier continuación probabilística, sea \(g(t)\in[0,1]\)
su probabilidad de éxito dada la terna o cuaterna \(t\). Por tanto
\[
 P_{\rm perm}(\text{éxito})
 =\frac{\sum_{t\text{ distinto}} g(t)}{(M)_r}
 \le \frac{M^r}{(M)_r}\,P_{\rm func}(\text{éxito}).
\]
La desigualdad vale también cuando la continuación conoce la salida y
parte del estado, siempre que pueda simularse tal como se definió en
la reducción PRP. En Lean, \proofat{PTG3.Reduction.retained_success_count}{retained\_\allowbreak{}success\_\allowbreak{}count} demuestra la
desigualdad de sumas y \proofat{PTG3.Reduction.paths_falling}{paths\_\allowbreak{}falling} cuenta las continuaciones
sin reemplazo. En un tramo real se paga además \(\alpha_4(T)\).
Sea \(L\le2N\) el máximo público de tramos en una época (se completan
los tramos ausentes con la identidad del experimento). Con
\(\gamma=M^4/(M)_4\), cada sustitución satisface
\(p_j\le\alpha_4(T)+\gamma p_{j+1}\). Para semilla uniforme,
\[
 p_0\le \gamma^L\bigl(L\alpha_4(T)+p_L\bigr).
\]
\proofat{PTG3.Reduction.success_hybrid}{success\_\allowbreak{}hybrid} prueba la inducción con denominadores
multiplicados en ambos miembros. Para las distribuciones finitas con monedas
aleatorias, todos los sucesos tienen probabilidades racionales.

Sea el objetivo una cadena de \(n\) posiciones de bit distintas
de salidas ocultas, prefijadas o elegidas sólo con información independiente
de esos bits en el experimento ideal. En el último experimento los bits
objetivo siguen uniformes dado el historial concedido, el estado terminal y
las entradas elegidas a partir de lo observable. La probabilidad de una
conjetura es \(p_L=2^{-n}\) aunque el objetivo abarque varias peticiones:
\proofat{PTG3.Reduction.uniform_guess_count}{uniform\_\allowbreak{}guess\_\allowbreak{}count} cuenta exactamente una cadena correcta por
cada vista y elección de monedas.
El historial concedido excluye copias de la cadena objetivo o funciones
laterales de ella que revelen sus bits. El estado anterior y el estado
posterior de la misma petición constituyen otra vista. Para una semilla a
distancia \(\eta_0\) del ideal en la vista previa, la cota final es
\[
 \Pr[\text{acierto}]
 \le \min\{1,\eta_0+\gamma^L(L\alpha_4(T)+2^{-n})\}.
\]
La cota de muestreo es multiplicativa para este suceso. La cota aditiva de
distinción de la sección~\ref{sec:ctr-prp} sigue aplicándose al transcript.

Otra vía para el mismo suceso usa una cota puntual de densidad de la
semilla clásica. Si para todo contexto previo admisible \(e\) y estado \(s\)
se cumple \(P(S_0=s\mid e)\le 2^d/2^{384}\), y la continuación condicionada
a \((s,e)\) es la misma en ambos modelos, basta multiplicar la cota de
semilla uniforme por \(2^d\). Se aplica \proofat{PTG3.Reduction.seed_density_transfer}{seed\_\allowbreak{}density\_\allowbreak{}transfer}
a los pesos de las continuaciones y se promedia sobre \(e\). El rótulo
\emph{full entropy} por sí solo no demuestra esa premisa sobre la vista.
Para valorar \(\min(240,n)\) bits de esfuerzo clásico (§214) hay que
justificar la calidad de la semilla y el valor de \(L\alpha_4(T)\) en el
presupuesto de trabajo que se alega. La eliminación del término aditivo de
cumpleaños resuelve una pérdida de la demostración. No asigna a AES una
ventaja numérica concreta \cite[§§121--122, 209 y 211--214]{ais}.

\paragraph{Cuánto acredita esa cota para el estado efectivo (.3).}
Hay un límite numérico incluso si la semilla es perfecta. Fijemos de antemano
una clave AES y tres textos distintos. Un test compara las tres respuestas del
oráculo con los tres cifrados calculados bajo la clave elegida. Para un oráculo
AES de clave uniforme, el test acepta con probabilidad al menos \(2^{-256}\):
la propia clave siempre pasa (\proofat{PTG3.Reduction.fixed_key_test_count}{fixed\_\allowbreak{}key\_\allowbreak{}test\_\allowbreak{}count}).
Para una permutación aleatoria, las tres imágenes fijadas son distintas y el
test acepta con probabilidad exactamente \(1/(M)_3\). Cuando \(T\) cubre
tres consultas, tres cifrados locales y las comparaciones, el test pertenece al
conjunto que define \(\alpha_4(T)\). En particular,
\[
 \alpha_4(T)\ge 2^{-256}-1/(M)_3,\qquad
 2048\bigl(2^{-256}-1/(M)_3\bigr)>2^{-246}.
\]
La segunda desigualdad se comprueba con enteros en
\proofat{PTG3.Reduction.segment_bound_floor_2048}{segment\_\allowbreak{}bound\_\allowbreak{}floor\_\allowbreak{}2048}. En una época de \(N=1024\) peticiones
con hasta \(L=2N\) tramos, el término \(\gamma^L L\alpha_4(T)\) del lado
derecho ya supera \(2^{-246}\). Incluso restringiendo toda la época a
\(L=1024\) tramos autónomos, la suma
\(1024(2^{-256}-1/(M)_3)+2^{-256}\) excede ese umbral
(\proofat{PTG3.Reduction.autonomous_bound_floor_1024}{autonomous\_\allowbreak{}bound\_\allowbreak{}floor\_\allowbreak{}1024}). Esta vía conservadora resulta
\emph{insuficiente} para establecer el criterio suficiente
\(p_{256}(T)\le2^{-246}\) que se propuso para .3 en dicho horizonte.
Es un límite de la cota superior por sustitución; la probabilidad real de
predicción puede ser mucho menor. Para justificar .3 se dispone del análisis
estructural de estados eficaces y su relación con §102, o se necesita una
reducción específica al predictor que evite pagar una ventaja PRP genérica
por cada tramo. Un conteo en el modelo de cifrado ideal puede apoyar esa
justificación si se identifica como tal.

\paragraph{2. Invertir una terna uniforme (.9).}
Con \(X=\bits^{256}\times\bits^{128}\) y \(Z=(\bits^{128})^3\), sea
\[
 F(k,v)=\bigl(\AES_k(\operatorname{inc}^{j}(v))\bigr)_{j=1}^3 .
\]
\(X\) y \(Z\) tienen exactamente \(2^{384}\) elementos
(\proofat{PTG3.Reduction.aes_domain_sizes}{aes\_\allowbreak{}domain\_\allowbreak{}sizes}). El codominio
incluye las ternas repetidas, de fibra vacía. Sea \(B\) un algoritmo que
devuelve un par candidato o el símbolo de fallo. Para \(x\sim U_X\) y
\(z\sim U_Z\) se obtiene
\[
 \Pr_z[B(z)\ne\bot\ \wedge\ F(B(z))=z]
       =\Pr_x[B(F(x))=x].
\]
Fijadas las monedas de \(B\), contamos los pares \((x,z)\) que verifican
\(F(x)=z\) y \(B(z)=x\). Cada objetivo invertido cuenta una sola vez incluso
si su fibra es grande. Los dos denominadores coinciden. Al promediar las
monedas se obtiene \proofat{PTG3.Reduction.randomized_double_count}{randomized\_\allowbreak{}double\_\allowbreak{}count}. La probabilidad de
encontrar \emph{alguna} preimagen al generar \(z=F(x)\) pondera las fibras
por su tamaño y exige el transporte de distribución del paso 3.

Un verificador elige \(v\) uniforme, obtiene del oráculo las tres palabras
de \(F(\cdot,v)\) y entrega únicamente esa terna a \(B\). Comprueba que el
contador candidato es \(v\) y coteja la clave candidata con el oráculo en
\(\operatorname{inc}^{j}(v)\), \(j=4,5,6\). Los seis puntos son distintos
incluso al envolver el contador (\proofat{PTG3.TransitionPreimage.six_points_fresh}{six\_\allowbreak{}points\_\allowbreak{}fresh}). Con el oráculo AES real, el par correcto
pasa siempre. En el mundo de permutación aleatoria, dados \(v\), las
primeras tres respuestas y las monedas del inversor, quedan \((M-3)_3\)
ternas equiprobables y como máximo una satisface las tres comprobaciones.
Así, para el inversor de terna uniforme,
\[
 \beta_3(T)\le \alpha_6(T')+\frac{1}{(M-3)_3}.
\]
\(T'\) incluye ejecutar el inversor y verificar el candidato. Los teoremas
\proofat{PTG3.Reduction.uniform_target_le_actual_verifier}{uniform\_\allowbreak{}target\_\allowbreak{}le\_\allowbreak{}actual\_\allowbreak{}verifier},
\proofat{PTG3.Reduction.paths_falling}{paths\_\allowbreak{}falling} y \proofat{PTG3.Reduction.candidate_at_most_one}{candidate\_\allowbreak{}at\_\allowbreak{}most\_\allowbreak{}one} comprueban
la primera comparación y los conteos de la segunda. La desigualdad restante
es la definición de ventaja PRP aplicada al verificador concreto.
La cota se deduce bajo la hipótesis PRP ordinaria para la clave uniforme.

\paragraph{3. Del estado terminal a la terna.}
Se fija la entrada de la transición que se intenta invertir, incluida la
ausencia de entrada; puede haber sido elegida en función del historial
permitido. Una preimagen candidata \(s^\star\) se comprueba ejecutando
\(\phi(s^\star,D)\). Si pasa, el preprocesado de esa misma petición
determina la clave de generación \(\widehat K^\star\) y el contador usado
para la salida \(w^\star\). El teorema \proofat{PTG3.TransitionPreimage.accept_iff_triple}{accept\_\allowbreak{}iff\_\allowbreak{}triple} enlaza
directamente \emph{nistGenerate} con \(F\): un estado candidato pasa el chequeo
si y sólo si \(F(\widehat K^\star,w^\star)\) coincide con la terna recibida
al deshacer el XOR de la entrada \(D\). \proofat{PTG3.TransitionPreimage.checker_equals_inversion}{checker\_\allowbreak{}equals\_\allowbreak{}inversion}
cubre el fallo explícito del inversor. \proofat{PTG3.TransitionPreimage.uniform_transition_le_aes_verifier}{uniform\_\allowbreak{}transition\_\allowbreak{}le\_\allowbreak{}aes\_\allowbreak{}verifier}
conecta después el conteo de preimágenes aceptadas del estado con el verificador
sobre seis consultas. Los tres teoremas de \texttt{TransitionPreimage.lean}
importan el modelo de 90A y la reducción finita de los pasos anteriores.

El chequeo de una preimagen usa siempre la función pública concreta
\(\phi\), también al evaluar el suceso de éxito en las distribuciones
sustituidas. Sea \(H\) todo el prefijo concedido, incluida la entrada real
elegida antes de esta petición. En el último tramo ya sustituido,
condicionados al prefijo completo (incluso a los valores internos del
prefijo), se consulta una función aleatoria nueva en cuatro puntos
distintos. Sus cuatro respuestas \((Y,C_1,C_2,C_3)\) son independientes
uniformes. Para todo \(h,y\) de probabilidad positiva y toda terna \(c\),
\[
 \Pr[C=c\mid H=h,Y=y]=M^{-3},\qquad
 \Pr[S^+=s\mid H=h,Y=y]=M^{-3}.
\]
La segunda igualdad resulta de que cada máscara fijada por \(h\) es una
biyección (\proofat{PTG3.Reduction.masked_target_fiber}{masked\_\allowbreak{}target\_\allowbreak{}fiber}). Por tanto el simulador puede
generar el prefijo y la última salida, añadir la entrada real, y entregar el
estado obtenido de la terna desafío mediante esa máscara. Esto incluye
salidas inmediatamente anteriores y, si se concede, la última salida. Las
dos igualdades pertenecen al modelo de funciones por tramo: volver a AES
cuesta las sustituciones de la ecuación siguiente. Se presupone compromiso
en el instante declarado, sin seleccionar la muestra mediante un predicado
sobre el estado terminal ni añadir telemetría dependiente de esa terna.
Cada preimagen aceptada proporciona un candidato para \(F\).
En consecuencia, para semilla a distancia \(\eta_0\), el éxito de
invertir en el perfil y vista definidos satisface
\[
 \Pr[\phi(s^\star,D)=s^+]
 \le \min\{1,\eta_0+\gamma^L[L\alpha_4(T)
                     +\alpha_6(T')+1/(M-3)_3]\}.
\]
La especialización a operación autónoma usa entrada nula. Si se admite que el inversor devuelva también una entrada elegida
\(D^\star\), la cota del último experimento se modifica. Para cada entrada
candidata fija \(d\), el objetivo uniforme \(z\) induce el estado terminal
\(z\oplus D(d)\), todavía uniforme e independiente del prefijo y de la
entrada real. Se cuentan sólo los éxitos en que el inversor devuelve
ese \(d\). La preimagen resultante invierte \(z\), por el mismo paso de
desenmascarado. Al sumar las clases de candidatos se sustituye
\(\beta_3\) por \(\min(1,(2^{128}+1)\beta_3)\) para el alfabeto del perfil
(entrada nula o una palabra). \proofat{PTG3.Reduction.chosen_input_union}{chosen\_\allowbreak{}input\_\allowbreak{}union} verifica esa
contabilidad, incluida la dependencia de la entrada elegida respecto del
desafío. Esta cota más débil conserva el ámbito del argumento cuando el
criterio de preimagen cambia y no implica una fuerza de 240 bits en esa
vista. El expediente identifica la interfaz al aplicar §§207 y 212--214. El paso de AES real a una permutación
sigue condicionado a la hipótesis de primitiva y a la simulación de la
sección~\ref{sec:ctr-prp}. Los teoremas de \texttt{Reduction.lean} comprueban
los conteos, el verificador y las desigualdades de composición. La
correspondencia de la vista con el TOE y la fuerza concreta de AES se
justifican por separado \cite[§§207, 212--214 y 1007]{ais}.
